\section{Discussion and Conclusion}
\label{sec:discussion}

Our work reframes the known fragility of sampling-based \ac{iDP}: the
excess vulnerability that prior work attributed to adversaries is what
rational play produces on its own. The Privacy Externality Game admits
no symmetric pure equilibrium among symmetric players: the diagonal
is a trough of each player's own utility, deviations from it are
self-reinforcing because budgets are strategic substitutes, and the
only symmetric equilibrium is a mixed lottery over who is exposed
(proven for two players at interior profiles, corroborated
numerically beyond). At every pure equilibrium
someone carries the redistributed risk, heterogeneous preferences sort
the exposed role onto whoever values privacy least, and the coalition
attack of \citet{kaiser2026your} is the coordinated limit of the same
incentive. The phenomenon is a property of the sampling-based
realisation---its sensitivity-based sibling, replayed through the
identical game, has a unique symmetric equilibrium---and its practical
bite is carried by attack power: at the measurement floor the
strategic pressure vanishes; against the strongest adversary the
deployed mechanism admits---the level any risk assessment must
price---it does not.

\paragraph{Limitations.} Four boundaries of the analysis are worth
stating plainly. First, the one-shot game assumes full information
and simultaneous play; our repeated Bayesian extension
(\cref{sec:brpeg}, \cref{app:bayesian}) covers a myopic follower and
a learning leader, but the fully two-sided repeated game---a target
who manages what their actions reveal~\cite{aumann1995repeated}---and
finite-horizon regret rates in the style of
\citet{balcan2015commitment} remain open. Second, our players are
silos choosing scalar budgets; per-record budget distributions within
a silo, and the granularity question they raise, are untreated.
Third, the continuous-action existence theorem covers the small-$p$
regime that \ac{DPSGD} occupies, not the saturation corner; and on coarse
grids the $n$-player symmetric game can pin fictitious play to the
diagonal (\cref{sec:H3}), so the $n \geq 3$ analogue of
\cref{prop:no-symmetric-ne} is conjectural. Fourth, our single
measured instance (\cref{sec:H1}) sits at the \ac{LiRA} measurement floor;
a measured instance with above-floor attack power is the missing
experiment between the analytical game and deployment reality.

\paragraph{Open questions.} Beyond closing those gaps: does the
$\Delta$-restricted game of \cref{prop:delta-neutralisation} admit a
unique (ideally symmetric) equilibrium, and at what model-utility
cost? Is there a budget-coordination mechanism that keeps
sampling-based \ac{iDP}'s utility advantage while provably removing the
externality---or is the sensitivity-based realisation's symmetric
equilibrium the price of strategic neutrality? And what should a
contributor-facing contract disclose: our aware-target result
(\cref{sec:H4}) says the coupling itself is the single most valuable
fact a contributor can know, which turns disclosure design into a
security control.

\paragraph{Conclusion.} Individual differential privacy promises that
each contributor controls their own risk. Under the sampling-based
realisation the promise fails not only against adversaries but at
equilibrium: risk is redistributed by rational play, invisibly to
every individual contract. In the hierarchy \eqref{eq:hierarchy}, the
contract is written at the mechanism-agnostic first level; the
redistribution takes place at the mechanism-specific second
level---the true bound on what any adversary can extract from the
deployed mechanism, which every player's choice moves and no
contract states; and any deployed attack realises some fraction of
that bound at the third. The remedy is not to police the
players but to complete the contract---bound the externality itself,
so that what each contributor is promised is what, whatever everyone
else does, they get.
